% =====================================================================
\section{Automated Hiring: Architecture, Failure Landscape, and Regulatory Trends}
\label{sec:landscape}
% =====================================================================

\subsection{Layered Architecture and Surfaces}

An evidence-based screening system spans four surfaces, shown in
Figure~\ref{fig:layers}. A candidate-facing portal collects consent, résumé, and
session media. An inference service hosts the perception and learned components.
A persistence layer event-sources the session so it can be reconstructed turn by
turn. A recruiter-facing application renders the audit. The separation matters
because the guarantees this article describes are enforced at boundaries between
these surfaces rather than inside any one of them.

\begin{figure}[t]
  \centering
  \includegraphics[width=\columnwidth]{fig05_layered_arch.pdf}
  \caption{Layered deployment architecture across the four surfaces.}
  \label{fig:layers}
\end{figure}

Session state is written to an append-only table with a sequence-assigning
database trigger. Question plans are deliberately \emph{not} persisted, so a plan
cannot be retrospectively edited to match an outcome --- a small design decision
that removes a whole category of post-hoc rationalisation.

\subsection{Failure Landscape}

Automated screening fails in ways that conventional information systems do not,
because its outputs are judgements about people and its inputs are difficult to
verify. Figure~\ref{fig:failures} organises the failure modes we encountered, and
Table~\ref{tab:threats} states their manifestations and defence implications.

\begin{figure*}[t]
  \centering
  \includegraphics[width=\textwidth]{fig06_failure_landscape.pdf}
  \caption{Failure modes in automated candidate screening. The four families are
  not independent: fabricated evidence typically originates in an unvalidated
  constant and is concealed by opaque aggregation.}
  \label{fig:failures}
\end{figure*}

\begin{table*}[t]
\caption{Failure modes in automated screening, their manifestations, and defence
implications.}
\label{tab:threats}
\small
\begin{tabular}{@{}p{2.6cm}p{6.2cm}p{6.6cm}@{}}
\toprule
Failure mode & Manifestation & Defence implication \\
\midrule
Fabricated evidence & A heuristic proxy is logged under the label of a
measurement; a confident pass is returned when the underlying record is absent &
Measurements must be structurally distinguishable from inferences; absence must
have its own representation \\
Silent component outage & A model fails to load, a fallback fills the gap, and
the status is never surfaced & Health of every perception component must be
rendered where a human will see it \\
Unvalidated constants & A similarity or timing threshold is chosen by reasoning
and never measured & Security-adjacent constants require calibration studies with
reported error rates \\
Opaque aggregation & A composite score is produced from hidden weights and cannot
be explained or contested & Aggregation should be omitted where the downstream
decision is made by a human anyway \\
Provenance gaps & Identity is verified once at login and assumed thereafter &
Provenance must be continuous, and unmeasured intervals must be visible \\
Instruction injection & Text in a candidate document is interpreted as a
directive by a downstream model & Planning must consume typed structures, never
raw document text \\
\bottomrule
\end{tabular}
\end{table*}

The first two rows are not hypothetical. Section~\ref{sec:deployment} reports an
audit of a reference implementation in which both occurred simultaneously: an
object detector that never successfully loaded, a luminance heuristic that filled
the gap and logged its output as camera-verified detections, and a status string
reporting readiness that was never rendered anywhere.

\subsection{Current State and Industry Trends}

Table~\ref{tab:trad-vs-current} contrasts traditional screening practice with
current automated practice and states what each advance costs.

\begin{table*}[t]
\caption{Comparison of traditional and current automated screening solutions.}
\label{tab:trad-vs-current}
\small
\begin{tabular}{@{}p{2.4cm}p{4.4cm}p{4.4cm}p{4.2cm}@{}}
\toprule
Aspect & Traditional & Current & Advances and limitations \\
\midrule
Output object & Interviewer notes and a rating & Composite score or ranking &
Higher throughput; loses the evidence that justified it \\
Identity & Visual confirmation in person & One-time check at login &
Scales; assumes continuity that remote settings do not guarantee \\
Process evidence & Observed directly by the interviewer & Captured and filed for
retrospective review & Available for fraud investigation; not fed back into the
live interview \\
Consistency & Varies by interviewer & Uniform by construction & Removes one bias
source; can entrench another at scale \\
Explanation & Reconstructed from memory & Feature attributions over a score &
Post-hoc; explains the model, not the decision \\
Contestability & Conversation with the recruiter & Appeal against an opaque
number & Formally available, practically difficult \\
Auditability & Paper trail, incomplete & Logged, but of the score not the
evidence & Satisfies retention; not sufficient for a bias audit \\
\bottomrule
\end{tabular}
\end{table*}

\begin{figure}[t]
  \centering
  \includegraphics[width=\columnwidth]{fig08_regulatory_timeline.pdf}
  \caption{Regulatory instruments governing automated employment decisions.
  Dates are statutory effective dates.}
  \label{fig:regulation}
\end{figure}

Regulation is converging on explainability. Figure~\ref{fig:regulation} places
the four instruments that bear most directly on this system. The Illinois
Artificial Intelligence Video Interview Act took effect in 2020 and requires
notice and consent for AI analysis of interview video. New York City Local Law
144 requires bias auditing and candidate notification for automated employment
decision tools. The European Union Artificial Intelligence Act classifies
employment-related systems as high risk, with corresponding transparency and
human-oversight obligations. The Illinois Human Rights Act as amended by
HB 3773 constrains AI use in employment decisions from 2026.

These instruments share a direction: an automated employment decision must be
explicable and contestable. A system architected around a scalar faces a retrofit
that its architecture resists, because the explanatory information was discarded
at design time. Figure~\ref{fig:standards} summarises the practices that follow.

\begin{figure*}[t]
  \centering
  \includegraphics[width=\textwidth]{fig09_standards.pdf}
  \caption{Compliance practices implied by the instruments in
  Figure~\ref{fig:regulation}.}
  \label{fig:standards}
\end{figure*}

\subsection{Landscape of Grounding, ML, and DL}

Figure~\ref{fig:taxonomy} organises the technologies examined in the remainder of
this article. The division is not arbitrary: grounding governs what may be
\emph{asserted}, ML measures properties of \emph{evidence}, and DL supplies
\emph{perception}. Assigning a task to the wrong family is itself a failure mode
--- using a learned model where a rule is correct produces an unauditable
decision, and using a rule where measurement is required produces an unvalidated
constant.

\begin{figure*}[t]
  \centering
  \includegraphics[width=\textwidth]{fig03_taxonomy.pdf}
  \caption{Taxonomy of technologies for evidence-based candidate screening.}
  \label{fig:taxonomy}
\end{figure*}

% =====================================================================
\section{Evidence Grounding as an Enabler}
\label{sec:grounding}
% =====================================================================

\subsection{Grounding-Centric Extraction}

The central mechanism of the system is a constraint on what a language model is
permitted to produce. Its rule: \textbf{a claim's text must be a verbatim
substring of the résumé.} The model's role is selection, not generation.

Naive substring matching is insufficient, and Figure~\ref{fig:gate} shows the
decision pipeline that replaces it. Beyond verbatim containment the gate
implements negation rejection, so that \emph{``I have not worked with
Kubernetes''} yields no Kubernetes claim despite containing the substring; hedge
rejection, so that \emph{``some exposure to Kafka''} is not a Kafka experience
claim; contrastive-conjunction splitting, so that \emph{``I used Django, but
never in production''} is split rather than truncated; and clause scoping, so
that matching occurs within clause boundaries rather than across them.

\begin{figure}[t]
  \centering
  \includegraphics[width=0.92\columnwidth]{fig10_grounding_gate.pdf}
  \caption{The grounding gate as a decision pipeline. Any branch marked
  \emph{no} discards the proposed claim; nothing downstream sees it.}
  \label{fig:gate}
\end{figure}

\begin{figure*}[t]
  \centering
  \includegraphics[width=\textwidth]{fig11_grounding_flow.pdf}
  \caption{End-to-end grounded extraction. Segmentation and normalisation remain
  deterministic; the model contributes span selection only.}
  \label{fig:groundflow}
\end{figure*}

\paragraph{Enforcement.}
The grounding module is forbidden from importing the AI module, and an
architecture test fails continuous integration if that import appears. The
guarantee is therefore a property of the build rather than of developer
discipline. Table~\ref{tab:grounding-adv} summarises what this buys.

\begin{table}[t]
\caption{Key advantages of evidence grounding.}
\label{tab:grounding-adv}
\small
\begin{tabular}{@{}p{2.3cm}p{5.1cm}@{}}
\toprule
Advantage & Basis \\
\midrule
Verifiable claims & Every claim is a substring with character offsets into the
source document \\
Injection resistance & Résumé text reaches the planner only as typed claims, so
an instruction in a document has no field to occupy \\
Auditable rejection & A discarded claim is discarded by a named rule, not by
model preference \\
Stable evaluation & The gate is deterministic, so its behaviour can be regression
tested at corpus scale \\
No fluency confound & Extraction quality is fidelity to the source, not the
plausibility of generated text \\
\bottomrule
\end{tabular}
\end{table}

\subsection{Integrity Functions Enabled by Grounding}

\begin{table*}[t]
\caption{Integrity functions enabled by grounding, and where each is exercised.}
\label{tab:grounding-features}
\small
\begin{tabular}{@{}p{3.0cm}p{5.6cm}p{6.4cm}@{}}
\toprule
Function & Mechanism & Application \\
\midrule
Claim provenance & Character offsets retained from extraction through to the
audit & A reviewer clicks a claim and lands on the candidate's own words \\
Prompt-injection defence & Planner consumes typed claims only & A résumé
containing an embedded instruction has no channel to the planner \\
Assertion-scope fidelity & Negation and hedge detection inside clause boundaries
& Denials and qualifications are not converted into claims \\
Quote fidelity in speech & Spoken-turn quotes must match the transcript verbatim,
including disfluency & An interviewer cannot attribute tidied speech to a
candidate \\
Non-fabrication of coverage & Unexamined claims carry an explicit verdict &
Absence of evidence is reported, not omitted \\
\bottomrule
\end{tabular}
\end{table*}

The last row deserves emphasis. The audit's verdict set is
\texttt{substantiated}, \texttt{notDemonstrated}, \texttt{contradicted}, and
\texttt{notExamined}, and the four are deliberately not orderable.
\texttt{notExamined} is load-bearing: without it an unprobed claim is simply
absent from the report, and absence reads as acceptance. Making non-measurement
explicit is what prevents an audit from being fabricated by omission.

\subsection{Measured Behaviour of the Gate}

The project's central claim is that the model selects but never authors. The
preceding subsections argue this structurally; this one measures it.

An offline deterministic harness drives the real pipeline modules read-only over
a corpus of 5{,}200 synthetic résumés spanning eight fields and four seniority
levels. No language-model key is required, so the result is reproducible by
anyone. Three quantities are measured across roughly 100{,}000 trials: the
verbatim true-positive rate, the paraphrase rejection rate, and the negation-trap
rejection rate. Figure~\ref{fig:gateresults} reports the outcome.

\begin{figure}[t]
  \centering
  \includegraphics[width=\columnwidth]{fig07_grounding_gate_results.pdf}
  \caption{Grounding-gate evaluation before and after the clause-splitter fix
  described below. Trial counts are shown at right.}
  \label{fig:gateresults}
\end{figure}

\paragraph{The defect this surfaced.}
The initial run showed a 0.82\% false-rejection rate: 425 genuine claims wrongly
discarded. Every one contained a period-bearing token --- \texttt{Node.js},
\texttt{React.js}, \texttt{asp.net}, \texttt{Python 3.9} --- and \emph{zero}
occurred without one. The clause splitter treated a period, question mark, or
exclamation mark as a sentence terminal unconditionally, so
\emph{``\ldots using JavaScript, Node.js.''} split into \emph{``\ldots Node''}
and \emph{``js.''}, leaving a one-clause claim spanning two clauses where the
locator could never find it. The claim was then \emph{silently} rejected: no
error, no diagnostic, just an absent claim.

The fix constrains a terminal to count only when followed by whitespace or
end-of-string. A real sentence terminal always is; the period between two word
characters never is. Crucially, keeping dotted tokens whole makes clauses
\emph{longer}, which widens negation scope rather than narrowing it, so the fix
cannot weaken the guard --- and the measured before/after confirms it, with both
rejection rates unchanged at exactly 100\%.

This is a failure mode worth naming: the bug was \emph{silent}. Nothing crashed
and no test failed; claims simply went missing. Only a corpus-scale evaluation
with an explicit true-positive metric could surface it, which is an argument for
measuring the properties you claim rather than only testing the paths you wrote.

\paragraph{Honest scope.}
The corpus is synthetic and produced by a single generator, so 100\% extraction
coverage is an artifact of that format and is \emph{not} claimed as a general
result. The defensible claim is narrower and concerns the safety mechanism: the
gate rejected 100\% of model-style paraphrases and negation traps across roughly
100{,}000 trials while admitting 100\% of verbatim claims.

\subsection{Open Challenges in Grounded Extraction}

\begin{table}[t]
\caption{Future directions for grounded extraction and their expected impact.}
\label{tab:grounding-future}
\small
\begin{tabular}{@{}p{2.2cm}p{5.2cm}@{}}
\toprule
Direction & Expected impact \\
\midrule
Compositional claims & Extracting technology, artefact, metric, and agency as
typed relations rather than as a single span \\
Multilingual scope & Negation and hedge detection currently assume English clause
structure \\
Discontiguous spans & Some true claims are expressible only across
non-contiguous text; a discontiguous grounding relation would admit them without
loosening the gate \\
Speech grounding & Applying the constraint to spoken turns requires treating the
transcript as the source of truth, with its own error rate \\
\bottomrule
\end{tabular}
\end{table}

The gate's principal cost is recall on claims a résumé expresses across a
sentence boundary. Tightening that without loosening the safety property is the
main open problem in this component.
